%%%%%%%%%%%%%%%%%%%%%%%%%%%%%%%%%%%%%%%%%%%%%%%%%%%%%%%%%%%%%%%%%%%%%%%%%%%%
% CAPÍTULO 15: RESULTADOS ESPERADOS
%%%%%%%%%%%%%%%%%%%%%%%%%%%%%%%%%%%%%%%%%%%%%%%%%%%%%%%%%%%%%%%%%%%%%%%%%%%%

\chapter{Resultados esperados}
\label{ch:resultados-esperados}

\section{Resultados tangibles}

Al finalizar el TFM se espera disponer de los resultados recogidos en la Tabla~\ref{tab:resultados-tangibles}.

\begin{longtable}{p{1.1cm} p{8.4cm} p{5.14cm}}
  \caption{Resultados tangibles esperados.}
  \label{tab:resultados-tangibles}\\
  \hline
  \textbf{ID} & \textbf{Resultado} & \textbf{Formato} \\
  \hline
  \endfirsthead
  \multicolumn{3}{c}{\tablename\ \thetable{} -- continuación}\\
  \hline
  \textbf{ID} & \textbf{Resultado} & \textbf{Formato} \\
  \hline
  \endhead
  R1 & Estado del arte consolidado sobre Gemelos Digitales, movilidad compartida, \textit{matching} y simulación & Capítulo de la memoria + matriz bibliográfica \\
  R2 & Especificación de requisitos y casos de uso de la plataforma & Capítulo de la memoria + tablas RF/RNF \\
  R3 & Arquitectura de referencia para una plataforma de carpooling universitario P2P apoyada por un Gemelo Digital & Capítulo de la memoria + diagramas UML \\
  R4 & Modelo digital explícito de usuarios, vehículos, viajes, solicitudes y rutas & Capítulo de la memoria + diagrama de clases \\
  R5 & Implementación funcional de la plataforma (\textit{frontend} + \textit{backend}) & Repositorio de código \\
  R6 & Módulo de \textit{matching} conductor-pasajero & Repositorio + documentación técnica \\
  R7 & Módulo de Gemelo Digital con sincronización por eventos & Repositorio + documentación técnica \\
  R8 & Módulo de simulación con al menos tres escenarios reproducibles & Repositorio + conjuntos de datos \\
  R9 & Resultados de evaluación cuantitativa del efecto del carpooling & Capítulo de la memoria + tablas comparativas \\
  R10 & Identificación de limitaciones y líneas de trabajo futuro & Capítulo de conclusiones \\
  \hline
\end{longtable}

\section{Resultados esperados por nivel de evaluación}

La Tabla~\ref{tab:resultados-nivel} relaciona cada nivel de evaluación con su resultado esperado.

\begin{table}[H]
  \centering
  \begin{tabularx}{\textwidth}{Y{0.4348} Y{1.5652}}
    \hline
    \textbf{Nivel} & \textbf{Resultado esperado} \\
    \hline
    Funcional & Todos los casos de uso del MVP funcionan sin errores críticos. \\
    Técnico & El \textit{matching} responde en menos de 2 segundos para 50 usuarios y la sincronización del DT se produce en tiempo casi real. \\
    Experimental & Se obtiene una reducción medible de los vehículos utilizados en al menos un escenario de carpooling respecto al escenario base. \\
    \hline
  \end{tabularx}
  \caption{Resultados esperados por nivel de evaluación.}
  \label{tab:resultados-nivel}
\end{table}

\section{Contribuciones esperadas}

\begin{itemize}
    \item \textbf{Contribución metodológica:} una forma de integrar Gemelo Digital, \textit{matching} y simulación en un contexto de carpooling universitario P2P.
    \item \textbf{Contribución técnica:} una arquitectura modular y un modelo digital reutilizable en contextos similares.
    \item \textbf{Contribución experimental:} evidencia cuantitativa sobre el efecto del carpooling bajo escenarios controlados.
    \item \textbf{Contribución académica:} posicionamiento del TFM frente a la literatura existente, identificando el hueco que cubre.
\end{itemize}

\section{Qué no se espera obtener}

Para gestionar expectativas y proteger el alcance, no se espera obtener:
\begin{itemize}
    \item Un algoritmo de \textit{matching} novedoso ni una contribución teórica a la optimización.
    \item Un Gemelo Digital de alta fidelidad con sincronización bidireccional en tiempo real.
    \item Una validación con usuarios reales de la UAL.
    \item Un despliegue productivo a escala universitaria.
    \item Predicciones del comportamiento real de la comunidad universitaria.
\end{itemize}
